% Guion completo de la sustentación (15 minutos)
% Compilar: tectonic guion.tex
\documentclass[11pt,a4paper]{article}
\usepackage[spanish,es-nodecimaldot]{babel}
\usepackage[margin=2.1cm]{geometry}
\usepackage{xcolor}
\usepackage{amsmath,amssymb}
\usepackage{booktabs}
\usepackage{enumitem}
\usepackage[most]{tcolorbox}
\usepackage[hidelinks]{hyperref}

\definecolor{AzulIcesi}{RGB}{0,51,102}
\definecolor{VerdeIcesi}{RGB}{0,122,77}
\definecolor{GrisTexto}{RGB}{90,90,90}
\definecolor{Naranja}{RGB}{190,90,0}

\setlength{\parindent}{0pt}
\setlength{\parskip}{5pt}
\linespread{1.08}

% Encabezados
\newcommand{\parte}[3]{%
  \clearpage
  {\color{AzulIcesi}\rule{\linewidth}{1.2pt}}\par\vspace{-2pt}
  {\Large\bfseries\color{AzulIcesi} Parte #1 \,·\, #2}\hfill{\large\bfseries #3}\par
  \vspace{2pt}}
\newcommand{\diapo}[3]{%
  \par\vspace{8pt}
  {\large\bfseries Diapositiva #1 \,·\, #2}\hfill{\normalsize\bfseries\color{AzulIcesi} #3}\par
  \vspace{1pt}}
% Lo que hay que señalar en pantalla
\newcommand{\senala}[1]{{\small\color{AzulIcesi}\textbf{Señalar:} #1}\par}
% Marca de tiempo al final de cada parte
\newcommand{\reloj}[1]{\par\vspace{4pt}{\small\color{GrisTexto}\textbf{Control de tiempo:} al terminar esta parte
  deberías ir en \textbf{#1}.}\par}
% Transición a la parte siguiente
\newcommand{\pasa}[1]{\par{\small\color{GrisTexto}\textbf{Transición:}} \emph{#1}\par}

% Letra pequeña: para entender (no se lee en voz alta)
\newtcolorbox{nota}[1][Para entender (no se lee)]{breakable, enhanced, colback=black!3, colframe=black!25,
  boxrule=0.3pt, arc=1pt, left=5pt, right=5pt, top=3pt, bottom=3pt, fontupper=\footnotesize,
  title={#1}, fonttitle=\footnotesize\bfseries, coltitle=GrisTexto, colbacktitle=black!7,
  attach title to upper={\ \ }, before skip=4pt, after skip=6pt}
% Párrafo opcional: se dice solo si vas bien de tiempo
\newtcolorbox{opcional}{breakable, enhanced, frame hidden, colback=white, borderline west={2pt}{0pt}{VerdeIcesi},
  left=8pt, right=0pt, top=1pt, bottom=1pt, before skip=4pt, after skip=4pt,
  before upper={{\color{VerdeIcesi}\scriptsize\bfseries OPCIONAL}\par\vspace{1pt}}}

\begin{document}

% ------------------------------------------------------------------
\begin{center}
{\LARGE\bfseries\color{AzulIcesi} Guion de la sustentación}\par\vspace{4pt}
{\large SEMP: detector de espira para reconocer un Spark y un Kwid}\par\vspace{2pt}
{\color{GrisTexto} Electromagnetismo · Universidad Icesi · 15 minutos · 25 diapositivas}
\end{center}

\section*{Cómo usar este guion}

El profesor escoge quién presenta, así que \textbf{todos} deben poder decir el guion completo. Cada diapositiva
trae tres tipos de texto:

\begin{itemize}[leftmargin=1.2em, itemsep=2pt]
 \item \textbf{Texto normal:} lo que se dice en voz alta. No hay que leerlo palabra por palabra, pero sí decir
       todas las ideas y todos los números en \textbf{negrita}.
 \item {\color{AzulIcesi}\textbf{Señalar:}} qué parte de la imagen mostrar con el puntero mientras se habla.
 \item \textbf{Recuadros grises en letra pequeña:} la explicación de lo que está detrás. No se leen; sirven para
       entender y para responder si el profesor pregunta.
\end{itemize}

Los párrafos con una barra verde y la marca {\color{VerdeIcesi}\scriptsize\bfseries OPCIONAL} se dicen solo si
vas bien de tiempo. Sin ellos la presentación dura unos 15 minutos hablando tranquilo; los opcionales suman unos
3 minutos más, para cuando vas rápido. Al final de cada parte hay
un \textbf{control de tiempo}: si vas más de 30 segundos atrasado, salta los opcionales de la parte siguiente.

\begin{center}
\begin{tabular}{@{}clcl@{}}
\toprule
Parte & Tema & Diapositivas & Tiempo \\
\midrule
1 & El problema y el esquema D.E.P.A. & 1 -- 4 & 0:00 -- 2:35 \\
2 & La física: por qué el carro cambia la espira & 5 -- 8 & 2:35 -- 5:30 \\
3 & El modelo: inductancia y oscilador LC & 9 -- 10 & 5:30 -- 7:25 \\
4 & Métodos numéricos: Euler y RK4 & 11 -- 14 & 7:25 -- 10:10 \\
5 & Simulación y demostración en vivo & 15 -- 19 & 10:10 -- 13:20 \\
6 & Resultados, limitaciones y cierre & 20 -- 25 & 13:20 -- 15:00 \\
\bottomrule
\end{tabular}
\end{center}

\begin{nota}[Ritmo]
Una persona que habla tranquila en una exposición dice unas 140 palabras por minuto. El texto normal de este guion
tiene unas 2\,100 palabras y los opcionales unas 400, así que el texto normal con la demostración da unos 15
minutos, y con todos los opcionales unos 18. Los tiempos de cada diapositiva se calcularon con ese ritmo. Ensayar una vez con cronómetro es lo que más ayuda: si una parte se pasa del tiempo, quitar opcionales;
si queda corta, decir los opcionales o explicar uno de los recuadros grises con sus palabras.
\end{nota}

% ==================================================================
\parte{1}{El problema y el esquema D.E.P.A.}{0:00 -- 2:35}

\diapo{1}{Portada}{0:00 -- 0:25}
\senala{el carro llegando a la talanquera.}

Buenos días, profesor; buenos días a todos. Somos Luis Ernesto Rodríguez, Sebastián Santander, Kevin Lozano,
Juan Pablo Corchuelo y Mateo López. Nuestro proyecto se llama \textbf{SEMP}, Sensores ElectroMagnéticos de
Proximidad, y lo que hicimos fue modelar y simular el detector que abre la talanquera de un parqueadero. La imagen
de la portada la generamos con inteligencia artificial y es solo ilustrativa.

\begin{nota}
Cada diapositiva que usa una imagen hecha con IA lo dice en una esquina. Decirlo desde el principio evita la
pregunta: la guía permite usar IA generativa siempre que se cite y se use con criterio.
\end{nota}

\diapo{2}{Contenido}{0:25 -- 0:40}
\senala{la captura del simulador: el Spark parado sobre la espira.}

Vamos en este orden: modelo teórico, implementación numérica, simulación, resultados y conclusiones. La imagen es
una captura de nuestro simulador: un Spark parado sobre la espira, con las líneas del campo magnético que salen del
piso. Más adelante la vamos a ver funcionando.

\diapo{3}{Introducción}{0:40 -- 1:30}
\senala{el collage de aplicaciones: fábrica, semáforo, parqueadero.}

Un sensor inductivo detecta metal sin tocarlo: en las fábricas cuenta piezas en una banda, en algunos semáforos
sabe si hay carros esperando y en los parqueaderos abre la talanquera.

\begin{opcional}
Casi nunca se ve, porque va enterrado. Si uno mira el pavimento justo antes de una talanquera, a veces se nota un
cuadrado cortado en el asfalto: ahí adentro está la espira de alambre.
\end{opcional}

Nuestra pregunta fue: ¿esa misma espira, además de abrir la talanquera, puede decir \textbf{qué carro llegó}? Sin
cámara y sin oprimir nada. El profesor nos propuso un caso concreto: distinguir un \textbf{Chevrolet Spark GT} de
un \textbf{Renault Kwid}, dos carros muy comunes en Colombia y casi del mismo tamaño.

Nos pusimos tres objetivos: modelar el detector con las leyes de \textbf{Faraday y Ampère}; resolver el modelo
analíticamente y con \textbf{Euler y Runge-Kutta 4} programados por nosotros; y armar un \textbf{simulador} para
medir si de verdad reconoce cada carro.

\begin{nota}
¿Por qué sin cámaras? La cámara necesita luz, se ensucia, se moja y es más costosa; la espira va protegida bajo el
asfalto y dura años. Esa comparación vuelve en la diapositiva 21, así que aquí basta con mencionarla.
\end{nota}

\diapo{4}{Esquema D.E.P.A.}{1:30 -- 2:35}
\senala{las cuatro cajas de izquierda a derecha y, al final, el recuadro verde de abajo.}

Para ordenar el proyecto usamos el esquema \textbf{D.E.P.A.} que nos dio el profesor. La idea es la misma del
espectrómetro de masas que vimos en clase: primero se divide el aparato en etapas, y cada etapa usa solo la física
que necesita.

A la izquierda están los insumos, los lugares donde se usan estos sensores, y después las cuatro letras.
\textbf{D, Detección:} qué detectamos; aquí son dos cosas, si hay carro y cuál es. \textbf{E, Emisiones:} cómo
responde el carro al campo; el chasis responde con corrientes de Foucault. \textbf{P, Propiedades} de la onda: la
frecuencia de \textbf{55 kilohertz}, la energía y la potencia. Y \textbf{A, Alerta:} lo que entrega el sensor; la
frecuencia sube, se abre la talanquera y el sistema dice el modelo.

Uniendo el contexto de inicio y el de impacto sale el caso de estudio que llevamos al simulador: una espira de
\textbf{1.8 metros} bajo un Spark o un Kwid.

\begin{opcional}
El contexto de inicio es un parqueadero que necesita saber qué carro entra, y el de impacto es abrir la talanquera
y dejar registrado qué modelo entró.
\end{opcional}

\begin{nota}
El profesor pidió decir \emph{alerta} y no \emph{señal}, para no confundirla con la señal eléctrica. Si pregunta
para qué sirve el esquema: delimita la matemática. Por eso en la diapositiva 8 aparecen solo dos ecuaciones de
Maxwell y no las cuatro. En el espectrómetro de masas las etapas eran acelerar (conservación de la energía),
desviar (fuerza magnética) y medir el radio; aquí son detectar, responder, medir la frecuencia y alertar.
\end{nota}

\pasa{``¿Y por qué un carro cambia algo en un alambre que está enterrado? Esa es la parte física.''}
\reloj{2:35}

% ==================================================================
\parte{2}{La física: por qué el carro cambia la espira}{2:35 -- 5:30}

\diapo{5}{Qué pasa físicamente}{2:35 -- 3:35}
\senala{las líneas de campo que suben desde la espira y la placa encima; seguir los cinco pasos.}

Esto es lo que pasa, en cinco pasos.

\textbf{Uno:} por la espira pasa una corriente alterna, que va y viene 55 mil veces por segundo, y crea un campo
magnético que sube desde el pavimento y cambia todo el tiempo.

\textbf{Dos:} cuando un carro se para encima, ese campo atraviesa el chasis de acero. Como el campo cambia, por la
\textbf{ley de Faraday} se induce un campo eléctrico, que en el metal mueve cargas en círculos: son las
\textbf{corrientes de Foucault}.

\textbf{Tres:} por la \textbf{ley de Lenz}, esas corrientes crean su propio campo, en contra del de la espira:
el chasis bloquea parte del flujo.

\textbf{Cuatro:} con menos flujo por cada amperio, la \textbf{inductancia de la espira baja}, y eso hace que la
frecuencia del detector \textbf{suba}.

\textbf{Cinco:} entre más cerca el chasis, más fuerte el efecto. Esa es la clave para distinguir los carros.

\begin{nota}
\textbf{Inductancia:} es el flujo magnético que atraviesa la espira por cada amperio de corriente, $L=\Phi/I$. Es
la misma idea de la capacitancia de clase ($C=Q/V$, carga por voltio), pero con flujo por amperio. Si el chasis le
quita flujo, $L$ baja.

\textbf{Imagen mental útil:} el chasis se comporta como un espejo. Es como si encima hubiera una ``espira imagen''
con la corriente al revés; esa espira imagen le resta flujo a la real. Entre más cerca el espejo, más resta.

\textbf{Si preguntan el sentido de las corrientes:} si en un instante la corriente de la espira va en sentido
horario visto desde arriba y está creciendo, el campo bajo la espira apunta hacia abajo y crece; en el chasis se
induce una corriente antihoraria, que crea un campo hacia arriba y se opone al cambio.
\end{nota}

\diapo{6}{Representación del sistema}{3:35 -- 4:10}
\senala{en la figura del corte: la espira en la ranura, el cable al detector y el chasis con las cotas de 20 y
23 cm.}

Esta figura es el corte del sistema. Abajo, en una ranura de \textbf{5 centímetros} en el pavimento, va la espira:
\textbf{tres vueltas} de alambre. Un cable la conecta al detector, que tiene adentro un \textbf{oscilador LC}. Y
arriba está el chasis del carro.

Fíjense en un detalle: la distancia que importa no es la altura del carro al piso, sino la distancia hasta la
espira, que está enterrada. Por eso son \textbf{20 centímetros para el Spark y 23 para el Kwid}: la altura libre
más los 5 de la ranura.

\begin{nota}
La figura se hizo con la skill Anatomy de Claude, y se cita en la diapositiva; las medidas que muestra salen del
modelo en Python. En la vida real la ranura se rellena con un sellante, y por eso la espira aguanta años de tráfico
sin mantenimiento.
\end{nota}

\diapo{7}{Caso de estudio: Spark frente a Kwid}{4:10 -- 4:40}
\senala{la tabla; primero las dos primeras filas y después la fila en negrita (15 cm frente a 18 cm).}

Estos son los dos carros. Miden casi lo mismo, así que por tamaño no se distinguen. La diferencia está en la fila
en negrita: la \textbf{altura libre al piso}, \textbf{15 centímetros el Spark y 18 el Kwid}.

\begin{opcional}
Miren las dos primeras filas: 3.64 por 1.60 metros el Spark y 3.68 por 1.59 el Kwid. El Kwid es más alto porque
está diseñado con un estilo más de camioneta.
\end{opcional}

Supusimos tres cosas: que el chasis es una placa plana de acero, que el carro tapa casi el 90~\% de la espira y
que el régimen es cuasi-estático, que explicamos en la siguiente.

\begin{opcional}
La espira es la estándar de instalación: un cuadrado de 1.8 metros, que son 6 por 6 pies, con tres vueltas de
alambre calibre 14, enterrada a 5 centímetros.
\end{opcional}

\begin{nota}
\textbf{$\eta$ (eta)} es la fracción del área de la espira que queda bajo el carro. El carro mide 1.60 m de ancho y la
espira 1.8 m, así que no la tapa completa: 0.89 con el Spark y 0.88 con el Kwid.

\textbf{El acero:} conductividad $\sigma = 6\times10^6$ S/m (unas diez veces menos que el cobre) y permeabilidad
relativa $\mu_r = 100$. \textbf{El circuito:} $R = 0.6~\Omega$ y $C = 0.1~\mu$F. Las medidas de los carros salen de
las fichas técnicas (diapositiva 23). Las fotos de los carros son genéricas, hechas con IA: no son los modelos
reales, y la diapositiva lo dice.
\end{nota}

\diapo{8}{Solo las ecuaciones de Maxwell necesarias}{4:40 -- 5:30}
\senala{la ley de Faraday, luego la de Ampère y al final el recuadro de la profundidad de penetración.}

Como pidió el profesor, usamos solo las ecuaciones de Maxwell que hacen falta. La central es la \textbf{ley de
Faraday}: un flujo que cambia induce un campo eléctrico. En el chasis eso da las corrientes de Foucault, y en la
espira, el voltaje $L\,di/dt$ del circuito. La segunda es la \textbf{ley de Ampère}, sin corriente de
desplazamiento, que dice cuánto campo crea la espira. Gauss magnética solo entra como condición de frontera.

Uno se puede preguntar: si el acero es magnético, ¿no debería subir la inductancia? La respuesta está en el
recuadro, la \textbf{profundidad de penetración}. A 55 kilohertz el campo solo entra \textbf{0.087 milímetros} en el
acero, menos que una hoja de papel. El chasis se porta como un conductor perfecto y \textbf{expulsa} el flujo.

\begin{nota}
\textbf{Cuasi-estático (por qué no hay corriente de desplazamiento):} la longitud de onda a 55 kHz es
$\lambda = c/f = 5.4$ km, unas 3\,000 veces el tamaño de la espira. Entonces el campo se acomoda prácticamente al
instante en toda la región y no hay onda que se propague. Además, dentro del acero la corriente de desplazamiento es
$5\times10^{-13}$ veces la de conducción.

\textbf{Profundidad de penetración:} $\delta = \sqrt{2/(\omega\mu_0\mu_r\sigma)}$. Entre más alta la frecuencia, más
conductor o más magnético el material, menos entra el campo. Una hoja de papel mide unos 0.1 mm.

\textbf{¿Y cuándo gana el magnetismo?} A frecuencias bajas. En el segundo escenario (diapositiva 19), a 1 kHz, una
pieza de acero hace \emph{subir} $L$ un 5.5~\% y una de cobre la hace bajar un 8.5~\%. Esa diferencia es la que
permite distinguir materiales.
\end{nota}

\pasa{``Ya sabemos qué pasa. Ahora le ponemos números.''}
\reloj{5:30}

% ==================================================================
\parte{3}{El modelo matemático}{5:30 -- 7:25}

\diapo{9}{Inductancia y efecto del chasis}{5:30 -- 6:20}
\senala{la fórmula de $L_0$, después la integral y al final la tabla con $-12.9$~\% y $-11.4$~\%.}

Primero, la espira sola. La aproximamos por una circular con la misma área, de radio 1.016 metros, y con la
fórmula de arriba, con tres vueltas, da \textbf{83 microhenrios}.

\begin{opcional}
Para comprobarlo usamos la regla que aplican los instaladores de espiras, que da unos 71 microhenrios. Es el mismo
orden de magnitud, así que el valor es razonable.
\end{opcional}

Para el chasis usamos la solución de \textbf{Dodd y Deeds}, un resultado clásico de 1968 que resuelve Maxwell para
una bobina encima de una placa conductora y da el cambio de impedancia, $\Delta Z$. Lo importante de la integral es
que la \textbf{parte imaginaria} de $\Delta Z$ cambia la inductancia y la \textbf{parte real}, la resistencia. Eso
se multiplica por $\eta$, la fracción de la espira que el carro tapa.

En la tabla: con el Spark la inductancia baja un \textbf{12.9~\%} y con el Kwid un \textbf{11.4~\%}. Baja más con
el Spark porque está más cerca.

\begin{opcional}
La resistencia sube un poquito: es la energía que se pierde calentando el chasis con las corrientes de Foucault.
\end{opcional}

\begin{nota}
\textbf{La fórmula de $L_0$:} $L_0=\mu_0N^2a\,[\ln(8a/r_w)-2]$, con $a$ el radio de la espira y $r_w = 0.8$ mm el
radio del alambre. Va con $N^2$ porque cada vuelta crea campo y además cada vuelta lo recoge.

\textbf{La integral, por partes:} suma la respuesta del chasis a todas las ``formas'' posibles del campo (cada
valor de $\alpha$). $J_1$, una función de Bessel, describe la forma de una espira circular. $\Gamma$ compara la
placa con el aire: si la placa fuera un conductor perfecto, $\Gamma \to -1$, que es justo la espira imagen de la
diapositiva 5. $F$ lleva el factor $e^{-\alpha d}$, que hace que el efecto caiga rápido con la distancia $d$. La
integral se calcula numéricamente con SciPy; el error de la malla es de $1.3\times10^{-6}$.

\textbf{De $\Delta Z$ a $L$ y $R$:} una bobina tiene impedancia $Z = R + j\omega L$. Por eso
$L_\text{eff} = L_0 + \operatorname{Im}\Delta Z/\omega$ y $R_\text{eff} = R_0 + \operatorname{Re}\Delta Z$. En
números: $\Delta L = -10.7~\mu$H (Spark) y $-9.5~\mu$H (Kwid); $\Delta R = +0.064~\Omega$ y $+0.054~\Omega$.
\end{nota}

\diapo{10}{El oscilador LC, como el RC de clase}{6:20 -- 7:25}
\senala{la ecuación de Kirchhoff, el recuadro de la fórmula exacta y la tabla de frecuencias.}

¿Cómo se da cuenta el detector de que la inductancia cambió? Conecta la espira a un capacitor y forma un
\textbf{oscilador LC}. Lo analizamos igual que el \textbf{circuito RC de clase}: ley de mallas de Kirchhoff, con
la misma convención de signos. La diferencia es que ahora está la bobina, que aporta $L\,di/dt$, y como la
corriente es la derivada de la carga, sale una ecuación de segundo orden: $L\ddot q + R\dot q + q/C = 0$.

\begin{opcional}
Es la misma ecuación de una masa con un resorte y un poco de fricción: la carga va y viene entre las placas del
capacitor, y se va apagando poco a poco por la resistencia.
\end{opcional}

Proponiendo $q = e^{st}$, como en el RC, sale la frecuencia de oscilación: \textbf{55.25 kilohertz sin carro}. El
factor de calidad es \textbf{48}, o sea que la resistencia casi no frena la oscilación.

Y aquí está la idea central: la frecuencia va como uno sobre la raíz de $L$. Si $L$ \textbf{baja}, la frecuencia
\textbf{sube}. Con el Spark encima llega a \textbf{59.18 kilohertz} y con el Kwid a \textbf{58.69}. Esa es la
alerta: el detector ve que la frecuencia subió y abre la talanquera, y según cuánto subió, sabe qué carro es.

\begin{opcional}
Ojo con el recuadro: es la fórmula \emph{exacta}. Hay una aproximación que se usa mucho y que nosotros también
usamos al principio; en la parte numérica les contamos por qué la cambiamos.
\end{opcional}

\begin{nota}
\textbf{Analogía masa-resorte:} $L$ hace de masa (es la ``inercia'' de la corriente), $1/C$ hace de constante del
resorte, $R$ de fricción y $q$ de posición. Por eso $\omega_0 = 1/\sqrt{LC}$, igual que $\omega = \sqrt{k/m}$.

\textbf{Las raíces:} con $q=e^{st}$ queda $Ls^2 + Rs + 1/C = 0$, y $s = -R/2L \pm j\omega_d$. Como
$R/2L \approx 3\,600~\text{s}^{-1}$ es mucho menor que $\omega_0 \approx 347\,000$ rad/s, el circuito es
subamortiguado y $\omega_d \approx \omega_0$. La parte real da el decaimiento: en unos 10 ciclos la amplitud cae a
la mitad. Un detector real le repone energía en cada ciclo para que no se apague; para el modelo basta el circuito
libre.

\textbf{Condición inicial:} capacitor cargado y corriente cero, $q(0)=q_0$, $i(0)=0$, igual que el capacitor
cargado que se descarga en el RC de clase.

\textbf{Factor de calidad:} $Q = \omega_0L_0/R = 48$. Un $Q$ alto significa que la oscilación dura muchos ciclos y
que la frecuencia queda bien definida, lo que facilita medirla.

\textbf{Fórmula exacta:} como $f\propto 1/\sqrt{L}$, $f/f_0 = (L/L_0)^{-1/2} = (1+\Delta L/L_0)^{-1/2}$. Con el
Spark: $(1-0.1287)^{-1/2} = 1.0713$, o sea $+7.13~\%$, y $55.25 \times 1.0713 \approx 59.18$ kHz.
\end{nota}

\pasa{``Esta ecuación también la resolvimos numéricamente, y ahí pasó algo interesante.''}
\reloj{7:25}

% ==================================================================
\parte{4}{Métodos numéricos: Euler y RK4}{7:25 -- 10:10}

\diapo{11}{Métodos de análisis numérico}{7:25 -- 8:15}
\senala{el sistema con $(q, i)$, la fórmula de Euler, los cuatro $k$ de RK4 y el recuadro del procedimiento.}

Los dos métodos los programamos a mano, sin los integradores de las librerías. Como trabajan con ecuaciones de
primer orden, escribimos el oscilador como un sistema de dos: la derivada de la carga es la corriente, y la de la
corriente sale de la ley de mallas.

\textbf{Euler} es el más sencillo: toma la pendiente en el punto donde está y avanza en línea recta un paso $h$.
\textbf{Runge-Kutta 4} es más cuidadoso: en cada paso calcula la pendiente cuatro veces, una al inicio, dos en la
mitad y una al final, y avanza con un promedio de las cuatro.

Integramos 30 períodos y medimos la frecuencia \textbf{contando los cruces por cero} de la carga, que es como lo
hace un detector real.

\begin{nota}
\textbf{Pesos 1-2-2-1 de RK4:} es la misma idea de la regla de Simpson para integrales; los puntos del medio pesan
el doble. \textbf{Justificación de los métodos:} la guía pide los dos; Euler sirve de referencia y RK4 es el
estándar de precisión para ecuaciones diferenciales ordinarias.

\textbf{Cruces por cero:} hay dos por período. El código toma los de bajada y, entre los dos pasos donde la carga
cambia de signo, interpola con una recta para encontrar el instante exacto. La frecuencia es el número de períodos
entre el primer y el último cruce dividido por el tiempo que pasó.
\end{nota}

\diapo{12}{Validación numérica: convergencia}{8:15 -- 8:50}
\senala{la gráfica log-log y las pendientes de las dos rectas; después la tabla.}

¿Cómo sabemos que el código está bien? Lo probamos con el circuito RL de la bobina, que tiene solución exacta, como
el RC de clase. La gráfica muestra el error contra el tamaño del paso en escala logarítmica, donde la pendiente es
el orden del método. Medimos \textbf{1.01 para Euler y 4.02 para RK4}: justo lo que dice la teoría.

En la práctica eso significa que si duplicamos los pasos, el error de Euler baja a la mitad y el de RK4 baja
\textbf{16 veces}.

\begin{opcional}
En la tabla se ve con la frecuencia del oscilador: con 40 pasos por período, RK4 ya tiene un error de 5 en un
millón.
\end{opcional}

\begin{nota}
\textbf{Por qué 16:} $2^4 = 16$. En la tabla, de 20 a 40 pasos el error de RK4 pasa de $8.1\times10^{-5}$ a
$5.1\times10^{-6}$: $8.1/0.51 \approx 16$.

\textbf{Por qué sirve el RL:} tiene solución exacta, $i(t) = (V/R)(1-e^{-t/\tau})$ con $\tau = L/R$, igual que la
carga del capacitor en el RC de clase (en un $\tau$ llega al 63~\%). Si el código reproduce una solución que se
conoce, se puede confiar en él para la que no se conoce.

\textbf{Por qué log-log:} si $\text{error} = C\,h^p$, entonces $\log(\text{error}) = \log C + p\log h$, que es una
recta de pendiente $p$.
\end{nota}

\diapo{13}{En el oscilador, Euler explota}{8:50 -- 9:25}
\senala{la gráfica de energía: la línea roja de Euler subiendo y la azul de RK4 bajando.}

Pero en el oscilador Euler \textbf{explota}. En la gráfica de energía, la línea azul es RK4: la energía baja poco a
poco, porque la resistencia la gasta. La roja es Euler: la energía \textbf{crece} hasta \textbf{cien mil millones
de veces} la inicial, y eso es imposible, porque el circuito no tiene de dónde sacarla.

La razón es que en un oscilador cada paso de Euler multiplica la amplitud por un número un poquito mayor que uno,
sin importar qué tan pequeño sea el paso. Achicar el paso solo retrasa la explosión.

\begin{nota}
\textbf{Intuición geométrica:} sin pérdidas, la oscilación dibuja un círculo en el plano carga-corriente. Euler
avanza por la tangente, y la tangente de un círculo siempre lleva hacia afuera: cada paso cae en un círculo un poco
más grande, y la solución se va en espiral.

\textbf{La cuenta:} el factor por paso es $|1+j\omega h| = \sqrt{1+(\omega h)^2}$. Con 40 pasos por período,
$\omega h = 2\pi/40 \approx 0.157$, y la energía se multiplica por $1+(\omega h)^2 \approx 1.025$ en cada paso. En
30 períodos hay 1\,200 pasos: $1.025^{1200} \approx e^{29}$. Quitando la pérdida real por la resistencia
($e^{-3.9}$) queda $\approx 10^{11}$.

\textbf{¿Y en el RL?} Allí Euler sí es estable si $\Delta t < 2\tau$, porque la solución decae sin oscilar.
\textbf{RK4:} deja un 0.0197 de la energía inicial, que es exactamente $e^{-Rt/L}$.
\end{nota}

\diapo{14}{La simulación corrigió una fórmula}{9:25 -- 10:10}
\senala{el 6.44~\% de arriba, el 7.130~\% de abajo y el recuadro de la lección.}

Y esta es la parte de la que estamos más orgullosos. Al principio usamos la aproximación que aparece en muchos
textos: el cambio de frecuencia es menos la mitad del cambio de inductancia. Con el Spark daba \textbf{6.44~\%}.
Pero la simulación con RK4 daba \textbf{7.13~\%}, igual con cualquier paso, así que no era un error numérico.

La simulación tenía razón: esa aproximación solo vale para cambios pequeños, de un 1 o 2~\%, y un carro cambia la
inductancia casi un 13~\%. La fórmula exacta, la de la diapositiva 10, da 7.13~\%. Así que \textbf{el método
numérico decidió qué fórmula usar}, y corregimos todo el proyecto.

\begin{opcional}
Para nosotros esa es la mejor prueba de por qué vale la pena validar con simulación: la fórmula aproximada parecía
correcta y daba un número razonable, y solo la comparación con RK4 mostró el error.
\end{opcional}

\begin{nota}
\textbf{De dónde sale la diferencia:} $(1+x)^{-1/2} \approx 1 - x/2 + 3x^2/8 - \dots$ La aproximación se queda con
los dos primeros términos. Con $x = -0.129$, el término $3x^2/8$ vale 0.0062, o sea 0.6 puntos: casi toda la
diferencia entre 6.44 y 7.13. Si $x$ fuera 0.01, ese término valdría 0.004~\% y la aproximación sería perfecta.

\textbf{Si preguntan de dónde salió la aproximación:} la propuso la IA y aparece en textos de detectores; está
documentado en el uso de IA generativa (docs/08) como una de las decisiones críticas del equipo.
\end{nota}

\pasa{``Todo esto lo juntamos en un simulador. Vamos a verlo.''}
\reloj{10:10}

% ==================================================================
\parte{5}{Simulación y demostración en vivo}{10:10 -- 13:20}

\diapo{15}{Escena 3D del parqueadero}{10:10 -- 10:25}
\senala{las dos capturas: el Spark con 59.18 kHz y el Kwid con 58.69 kHz.}

El simulador está hecho en Python, con una escena 3D del parqueadero. Los carros tienen sus medidas reales, y todo
número que aparece en pantalla sale del modelo que acabamos de explicar.

\begin{opcional}
Usamos NumPy y SciPy para el modelo, SymPy para las deducciones, Streamlit para la interfaz y Three.js para la
escena.
\end{opcional}

\subsection*{Demostración en vivo \hfill {\normalsize\color{AzulIcesi} 10:25 -- 12:00}}

Cambiar al simulador, pestaña \textbf{Caso: vehículo}, con el Spark y la espira de 1.8 m.

\begin{enumerate}[leftmargin=1.4em, itemsep=3pt]
 \item \textbf{Dejar que llegue el Spark.} ``Aquí llega el Spark y se detiene sobre la espira. Miren el indicador:
       la frecuencia pasó de 55.25 a \textbf{59.18 kilohertz}, la pantalla dice SPARK y la talanquera se abre.''
 \item \textbf{Cámara ``A ras''.} ``Si bajamos la cámara al nivel del piso, se ve cómo las líneas de campo se
       acuestan debajo del chasis. Es el chasis expulsando el flujo, lo que explicamos con la profundidad de
       penetración.''
 \item \textbf{Cambiar al Kwid.} ``Ahora el Kwid. Como está tres centímetros más alto, la frecuencia sube menos:
       \textbf{58.69}. Y la pantalla dice KWID.''
 \item {\color{VerdeIcesi}\scriptsize\bfseries OPCIONAL} \textbf{Pestaña del oscilador LC.} ``Aquí se ve lo de
       Euler: con el mismo paso, RK4 oscila bien y Euler se sale de la gráfica.''
\end{enumerate}

\begin{nota}[Plan B]
Si el computador falla o la escena 3D no carga, no perder tiempo: volver a la diapositiva 15 y decir lo mismo
señalando las capturas, que muestran exactamente esos dos momentos. La diapositiva 18 tiene la vista a ras del piso.
\end{nota}

\diapo{16}{Lo que mide el detector}{12:00 -- 12:20}
\senala{la curva que baja con la altura y los dos puntos marcados.}

Esta gráfica resume el proyecto: entre más bajo el carro, más baja la inductancia, unos \textbf{0.46 puntos por
centímetro}. El detector hace el camino al revés: mide la frecuencia, saca la inductancia y con la curva estima la
altura. Da
\textbf{14.9 centímetros para el Spark y 18.0 para el Kwid}: los reconoce.

\begin{nota}
Tres centímetros de diferencia por 0.46 da unos 1.4 puntos, que es la diferencia entre $-12.87$ y $-11.39~\%$. La
curva siempre baja (es monótona), y por eso se puede invertir: a cada frecuencia le corresponde una sola altura.
\end{nota}

\diapo{17}{El carro cruzando la espira}{12:20 -- 12:40}
\senala{el trapecio: la subida, la parte plana y la bajada.}

Esto es lo que ve el detector cuando el carro pasa sin detenerse. La frecuencia forma un \textbf{trapecio}: sube
mientras el carro va tapando la espira, se queda arriba mientras la tapa completa y baja cuando sale. A 10 km/h dura
unos \textbf{2 segundos}, y el detector necesita solo \textbf{18 milisegundos} para medir: hay tiempo de sobra.

\begin{nota}
\textbf{2 segundos:} el carro tiene que recorrer su largo más el lado de la espira, $3.64 + 1.8 \approx 5.4$ m, a
2.8 m/s. \textbf{18 ms:} son 1\,000 ciclos de 55 kHz; contar 1\,000 ciclos da una resolución de 1/1\,000, o sea
0.1~\%. A 40 km/h la señal dura medio segundo: miles de mediciones por carro.
\end{nota}

\diapo{18}{Escenarios de diseño}{12:40 -- 12:55}
\senala{la columna ``Diferencia'' de las dos tablas y la vista a ras del piso.}

También usamos el simulador para diseñar. La mejor diferencia entre los dos carros sale con una espira de
\textbf{1.5 a 1.8 metros}, o sea \textbf{del tamaño del carro}. Y entre menos profunda, más señal.

\begin{opcional}
Si la espira es muy pequeña, el campo no alcanza a llegar al chasis; si es muy grande, el carro tapa solo una
parte.
\end{opcional}

\begin{nota}
El campo de una espira se debilita en una distancia parecida a su tamaño; una espira de 1 m ``ve'' poco a 20 cm de
altura comparada con una de 1.8 m. \textbf{Las vueltas:} $L$ crece con $N^2$ y $\Delta L$ también, así que el
cambio relativo no depende de $N$; las vueltas solo cambian la frecuencia de base.
\end{nota}

\diapo{19}{Segundo escenario: identificar el material}{12:55 -- 13:20}
\senala{el plano de impedancia, con cada material en una dirección distinta.}

El mismo modelo sirve para el otro caso del que habló el profesor, una línea de autopartes. Con una bobina pequeña
se mide el cambio de impedancia de cada pieza: la \textbf{fase} dice de qué metal es y la \textbf{magnitud}, a qué
distancia está. En las pruebas acertó el \textbf{100~\%}.

\begin{nota}
\textbf{Plano de impedancia:} se grafica $\Delta R$ contra $\Delta X$. Cada material sale en una dirección (un
ángulo) que depende de $\sigma$ y $\mu$ y casi no de la distancia; la distancia solo cambia qué tan lejos del origen
queda el punto. Cobre y aluminio bajan $L$; el acero, a 1 kHz, la sube. El 100~\% se logró con dos frecuencias y
16 promedios por medición. Se usan dos frecuencias porque cerca de 20 kHz el acero queda ``ciego'': el efecto
magnético y el de las corrientes se cancelan.
\end{nota}

\pasa{``Para cerrar, qué tan bien funciona y qué le falta.''}
\reloj{13:20}

% ==================================================================
\parte{6}{Resultados, limitaciones y cierre}{13:20 -- 15:00}

\diapo{20}{Resultados principales}{13:20 -- 13:40}
\senala{la tabla del caso del carro y después las cuatro viñetas.}

Resumiendo: la espira tiene \textbf{83 microhenrios} y oscila a \textbf{55.25 kilohertz}; RK4 da la frecuencia con
un error de 5 en un millón, mientras que Euler diverge; y en la simulación los dos carros se distinguen, con una
diferencia de \textbf{0.49 kilohertz}, que es fácil de medir.

\begin{opcional}
Y la potencia que el carro absorbe es de unos 80 microvatios con el Spark: el detector consume menos de un
milivatio en la espira.
\end{opcional}

\diapo{21}{Discusión}{13:40 -- 14:15}
\senala{las limitaciones en rojo y la tabla de comparación.}

El modelo coincide con cómo funcionan los detectores reales: el chasis baja la inductancia y el carro más bajo la
baja más.

\begin{opcional}
Comparada con una cámara o un sensor de ultrasonido, la espira funciona de noche y con lluvia y no tiene partes
expuestas; lo que no puede es leer el modelo exacto, como haría una cámara.
\end{opcional}

También tiene limitaciones, y queremos decirlas. Tratamos el chasis como una placa plana, así que el efecto que
calculamos es un \textbf{máximo}: el chasis real tiene huecos y partes a distintas alturas. Además, un Kwid muy
cargado baja a unos 16 centímetros y se podría confundir con un Spark. Eso se arregla con una segunda espira que
mida el largo del carro, y un análisis por elementos finitos daría el valor exacto.

\begin{nota}
Decir las limitaciones \emph{antes} de que pregunten muestra que se entiende el modelo. \textbf{Potencia:}
$P = \tfrac12 I^2\,\Delta R = \tfrac12(0.05)^2(0.064) = 80~\mu$W con 50 mA. \textbf{Elementos finitos:} dividir el
chasis real en millones de pedacitos y resolver Maxwell en cada uno con un programa como COMSOL; da el valor exacto
pero no se puede hacer a mano.
\end{nota}

\diapo{22}{Conclusiones}{14:15 -- 14:50}
\senala{los cinco puntos, uno por uno.}

Para cerrar, cinco conclusiones. \textbf{Uno:} con solo Faraday y Ampère se explica el detector. \textbf{Dos:} el
oscilador se resuelve con Kirchhoff, igual que el RC de clase. \textbf{Tres:} el detector reconoce los carros por su
altura: el \textbf{Spark a 59.18} y el \textbf{Kwid a 58.69 kilohertz}. \textbf{Cuatro:} RK4 converge con orden 4,
Euler explota en el oscilador, y la simulación nos corrigió una fórmula. \textbf{Cinco}, como trabajo futuro: el
modelo con el chasis real, la segunda espira y probarlo con un detector comercial.

\diapo{23 y 24}{Referencias y material complementario}{pasar rápido}

Las referencias y el enlace al repositorio, con el código, las 13 pruebas automáticas y el informe, quedan en estas
dos diapositivas.

\diapo{25}{Gracias}{14:50 -- 15:00}

Gracias. Quedamos atentos a sus preguntas.

% ==================================================================
\clearpage
\section*{Números que hay que saber de memoria}

\begin{center}
\small\begin{tabular}{@{}p{8.2cm}p{7.8cm}@{}}
\toprule
Qué & Valor \\
\midrule
Espira & 1.8 m, 3 vueltas, 5 cm de profundidad, $R = 0.6~\Omega$, $C = 0.1~\mu$F \\
Inductancia sin carro & 83 µH \\
Frecuencia sin carro / factor de calidad & 55.25 kHz / $Q = 48$ \\
Altura del Spark / Kwid (distancia a la espira) & 15 cm / 18 cm (20 cm / 23 cm) \\
Cambio de inductancia Spark / Kwid & $-12.9~\%$ / $-11.4~\%$ \\
Frecuencia con Spark / Kwid & \textbf{59.18 kHz / 58.69 kHz} \\
Diferencia entre los carros & 0.49 kHz \\
Altura estimada & 14.9 cm / 18.0 cm \\
Profundidad de penetración en el acero a 55 kHz & 0.087 mm \\
Orden de Euler / RK4 & 1 / 4 (medidos 1.01 / 4.02) \\
Aproximación lineal frente a fórmula exacta & 6.44~\% frente a 7.13~\% \\
Crecimiento de la energía con Euler & $\times 10^{11}$ en 30 períodos \\
Longitud de onda a 55 kHz & 5.4 km \\
\bottomrule
\end{tabular}
\end{center}

\section*{Preguntas más probables}

\sloppy\begin{description}[leftmargin=0pt, itemsep=4pt, font=\normalfont\bfseries, style=unboxed]
 \item[¿Qué ecuación de Maxwell es la clave?] Faraday: el campo que cambia induce las corrientes en el chasis y la
       fem en la espira.
 \item[¿Por qué baja la inductancia si el acero es magnético?] A 55 kHz el campo entra solo 0.087 mm y domina el
       apantallamiento de las corrientes de Foucault. El magnetismo gana a frecuencias bajas, como a 1 kHz en el
       escenario de las autopartes.
 \item[¿Cómo se relaciona con la clase?] El oscilador se resuelve con la ley de mallas y la misma convención de
       signos del RC, y la inductancia se define como la capacitancia, pero con flujo por amperio.
 \item[¿Por qué no usaron $\Delta f/f = -\tfrac12\Delta L/L$?] Porque solo vale para cambios pequeños; con un 13~\%
       de cambio la fórmula exacta da 7.13~\% y la aproximada 6.44~\%, y la simulación con RK4 confirmó la exacta.
 \item[¿Por qué Euler explota?] En un oscilador, cada paso de Euler multiplica la amplitud por $\sqrt{1+(\omega h)^2}$, que es
       mayor que uno: la solución sale en espiral.
 \item[¿Cómo saben que el código está bien?] Se probó con el circuito RL, que tiene solución exacta, y dio los
       órdenes teóricos (1.01 y 4.02). Además hay 13 pruebas automáticas.
 \item[¿Qué pasa si el Kwid va cargado?] Baja a unos 16 cm y se puede confundir con el Spark; se arregla midiendo
       también el largo con una segunda espira.
 \item[¿Por qué no una cámara?] La espira no le teme a la lluvia ni a la noche y no tiene partes expuestas; la cámara
       reconoce el modelo exacto, pero es más frágil y costosa.
 \item[¿El campo es peligroso?] No: unos 0.09 µT, cientos de veces menos que el campo magnético de la Tierra.
 \item[¿Dónde usaron inteligencia artificial?] Para escribir parte del código y los borradores, la figura del corte y
       las imágenes ilustrativas. Todo resultado se comprobó ejecutando el código, y la fórmula aproximada que propuso
       se corrigió con la simulación.
\end{description}

Más preguntas y respuestas en \texttt{docs/07-sustentacion.md}.

\section*{Antes de empezar}

\begin{itemize}[leftmargin=1.2em, itemsep=2pt]
 \item[$\square$] Simulador abierto en la pestaña \textbf{Caso: vehículo}, con el Spark y la espira de 1.8 m.
 \item[$\square$] Escena 3D probada en el computador de la exposición (necesita WebGL, no necesita internet).
 \item[$\square$] PDF de la presentación abierto como respaldo.
 \item[$\square$] Cada integrante ensayó el guion completo una vez con cronómetro.
\end{itemize}

\end{document}
